% concatenated from AndBach39.tex
\documentclass[12pt,reqno]{amsart}
\usepackage{fullpage}
\usepackage[top=1.5in, bottom=2in, left=2.5in, right=2.5in]{geometry}
\oddsidemargin = 0.5in \evensidemargin = 0.5in \textwidth = 5.5in
\usepackage{graphicx,amsmath,amssymb,amsfonts,amsthm,enumitem,bm,xcolor}
\usepackage{fullpage}
\usepackage{cleveref}


\usepackage[normalem]{ulem}

\usepackage{url}

%\usepackage[maxbibnames=10%backend=biber,
%%style=alphabetic,sorting=ynt
%]
%{biblatex}
%\addbibresource{unimodal.bib}

\usepackage{color}

%\usepackage{refcheck}
\DeclareMathOperator{\arccosh}{arccosh}

\newcommand{\N}{\mathbb{N}}
\newcommand{\R}{\mathbb{R}}
\newcommand{\CB}{\mathcal{B}}
\newcommand{\CM}{\mathcal{M}}
\newcommand{\CC}{\mathcal{C}}
\renewcommand{\Im}{\operatorname{Im}}
\newcommand{\df}{\dfrac}
\newcommand{\Arg}{\operatorname{Arg}}
\newcommand{\mult}{\operatorname{mult}}
\newcommand{\even}{\operatorname{even}}
\newcommand{\Gal}{\operatorname{Gal}}
\newcommand{\Res}{\operatorname{Res}}
\renewcommand{\Re}{\operatorname{Re}}
\renewcommand{\Im}{\operatorname{Im}}
\newcommand{\s}{{\sigma}}
\newcommand{\Sig}{{\Sigma}}
\newcommand{\sss}{\sum^{\infty}_{\substack{
			c,d=-\infty\\(c,d)=1\\c,d\text{ odd }}}}
\renewcommand{\arctan}{\tan^{-1}}
\renewcommand{\a}{\alpha}
\renewcommand{\b}{\beta}
\newcommand{\e}{\epsilon}
\renewcommand{\d}{{\delta}}
\newcommand{\g}{\gamma}
\newcommand{\G}{\Gamma}
\renewcommand{\l}{\lambda}
\renewcommand{\L}{\Lambda}
\renewcommand{\k}{\kappa}
\renewcommand{\o}{{\omega}}

\renewcommand{\O}{{\Omega}}
\renewcommand{\r}{{\rho}}
\renewcommand{\t}{\tau}
\newcommand{\x}{\xi}
\newcommand{\z}{\zeta}
\newcommand{\CCal}{{\mathcal{C}}}
\newcommand{\ECal}{{\mathcal{E}}}
\newcommand{\HCal}{{\mathcal{H}}}
\newcommand{\JCal}{{\mathcal{J}}}
\newcommand{\ZBbb}{{\mathbf{Z}}}
\newcommand{\qu}{{\underline{q}}}
\newcommand{\thre}{{\sqrt{3}}}
\newcommand{\qq}{{\(\dfrac q{\qu}\)^2}}
\renewcommand{\(}{\left\(}
\renewcommand{\)}{\right\)}
\renewcommand{\i}{\infty}
\newcommand{\Mod}[1]{\ (\mathrm{mod}\ #1)}
\newcommand{\os}[2]{\overset{#1}{#2}}
\newcommand{\ub}[1]{\underbrace{#1}}
\newcommand{\pa}[2]{\left(\frac{#1}{#2}\right)}
\newcommand{\pnoa}[2]{(\frac{#1}{#2})}
\newcommand{\flo}[1]{\lfloor #1\rfloor}
\newcommand{\Flo}[1]{\left\lfloor #1\right\rfloor}
\newcommand{\cei}[1]{\lceil #1\rceil}
\newcommand{\Cei}[1]{\left\lceil #1\right\rceil}
\newcommand{\obinom}{\genfrac{}{}{0pt}{}}
\newcommand{\od}{\mathrm{od}}
%\renewcommand{\b}{\mathbf{B}}
\numberwithin{equation}{section}
\theoremstyle{plain}
\newtheorem{theorem}{Theorem}[section]
\newtheorem{lemma}[theorem]{Lemma}
\newtheorem{notation}[theorem]{Notation}
\newtheorem{refproof}{Proof}
\newtheorem{note}[theorem]{Note}
\newtheorem{case}[theorem]{Case}

\newtheorem{problem}[theorem]{Problem}
\newtheorem{conjecture}[theorem]{Conjecture}
\newtheorem*{conjecture*}{Conjecture}
\newtheorem{corollary}[theorem]{Corollary}
\newtheorem{proposition}[theorem]{Proposition}
\newtheorem{example}[]{Example}

\theoremstyle{definition}
\newtheorem*{definition}{Definition}
\newtheorem*{remark}{Remark}

\newcommand\mycom[2]{\genfrac{}{}{0pt}{}{#1}{#2}}
\numberwithin{equation}{section}
\newcommand{\abs}[1]{\lvert#1\rvert}
\newcommand{\blankbox}[2]{%
	\parbox{\columnwidth}{\centering
		\setlength{\fboxsep}{0pt}%
		\fbox{\raisebox{0pt}[#2]{\hspace{#1}}}%
	}%
}

%Custom commands for Kathrin
\newcommand{\kbf}{\bf\color{green!40!black!60!}}
\newcommand{\cl}{\noindent\ul{Claim}: }
\renewcommand{\binom}[2]{\left(\begin{smallmatrix}#1\\\\#2\end{smallmatrix}\right)}
\newcommand{\smallbinom}[2]{(\begin{smallmatrix}#1\\#2\end{smallmatrix})}
\renewcommand{\pmod}[1]{\ \left( \mathrm{mod} \, #1 \right)}
\newcommand{\Pmod}[1]{\ ( \mathrm{mod} \, #1 )}

\setlist[enumerate]{leftmargin=*,label=\rm{(\arabic*)}}

\makeatletter
\@namedef{subjclassname@2020}{%
	\textup{2020} Mathematics Subject Classification}
\makeatother

\allowdisplaybreaks

%\makeatletter
%\newcommand{\vast}{\bBigg@{2}}
%\newcommand{\Vast}{\bBigg@{5}}
%\makeatother


\title{On congruence conjectures of Andrews and Bachraoui}
\author{Koustav Banerjee}
\author{Kathrin Bringmann}
\address{University of Cologne, Department of Mathematics and Computer Science, Weyertal 86-90, 50931 Cologne, Germany}
\email{kbanerj1@uni-koeln.de}
%\author{Kathrin Bringmann}
\email{kbringma@uni-koeln.de}
\author{Mohamed El Bachraoui}
\address{Dept. Math. Sci, United Arab Emirates University, PO Box 15551, Al-Ain, UAE}
\email{melbachraoui@uaeu.ac.ae}

\makeatletter
\@namedef{subjclassname@2020}{%
	\textup{2020} Mathematics Subject Classification}
\makeatother

\subjclass[2020]{05A17, 11P81, 11P83}
\keywords{congruences, mock theta functions, $q$-series, theta functions}



\begin{document}
	


\begin{abstract}

Andrews and the third author recently studied congruences for certain restricted two-color partitions.  They made two conjectures for Ramanujan-type congruences and a vanishing identity for the limiting sequence. In this paper, we settle these conjectures by relating the corresponding generating function to modular forms and mock theta functions. %and subsequently, using dissections of associated mock theta functions and eta quotients.


\end{abstract}

\maketitle


%{\bf KB: Details file needs to be updated. \color{blue} K: Updated the details file.}


%{\bf \color{blue} K: all the equation numbers are used, all references are used in the paper.}


%{\kbf To Study: 
%\begin{enumerate}
%	\item Whether $f(\z q,q^2;q^2)$ is a holomorphic part of a harmonic Maass form? 
%	\item Same question for $\nu(iq,i\left(\z q\right)^{\frac 12};q)$.
%\end{enumerate}}

%{\bf \color{blue} K: deleted all unused equation numbers.}
\section{Introduction and Statements of results} 

 In \cite[(1.1)]{AB}, Andrews and the third author considered restricted two-color partitions (see \cite[Definition 1]{AB}) $C(k,n)$ (for $k\in \mathbb{N}$), which count the number of two-color partitions of $n$, where the smallest part is odd and occurs at least once in blue color, every even part in blue color is at least $2k-1$ greater than the smallest odd part, and even parts of same color are distinct. They showed that
\[
C_k(q):=\sum_{n\ge 0}C\left(k,n\right)q^n=\sum_{n\ge 0}\frac{\left(-q^{2n+2k},-q^{2n+2};q^2\right)_{\infty}q^{2n+1}}{\left(q^{2n+1};q^2\right)^2_{\infty}},
\]  
where $(a)_n=(a;q)_n:=\prod_{j=0}^{n-1}(1-aq^j)$ ($n\in \mathbb{N}_0\cup \{\infty\}$). Letting $k\to \infty$, in \cite[Section 10]{AB} they introduced
\[
C(q):=\lim_{k\to \infty}C_k(q)=\sum_{n\ge 0}\frac{\left(-q^{2n+2};q^2\right)_{\infty}q^{2n+1}}{\left(q^{2n+1};q^2\right)^2_{\infty}}=:\sum_{n\ge 0}c(n)q^n.
\]
In \cite[Remark 3]{AB}, the authors then represented $C(q)$ as 
\begin{equation}\label{C}
C(q)=qA(q)S(q),
\end{equation}
where\footnote{The definition of $A(q)$ slightly differs (by a prefactor $q$) from Andrews and Bachraoui's definition of $A(q)$, see \cite[Remark 3]{AB}.}  %{\bf \color{blue} K: make a footnote on our notation and Andrews' notation.}
\begin{align}\label{A}
A(q)&:=\!\frac{\left(-q^2;q^2\right)_{\infty}}{\left(q;q^2\right)^2_{\infty}}\!=:\!\sum_{n\ge 0}a(n)q^n,\\\label{S}
 S(q)&:=\!\sum_{n\ge 0}\frac{\left(q;q^2\right)^2_n q^{2n}}{\left(-q^2;q^2\right)_n}\!=:\!\sum_{n\ge 0}s(n)q^n.
\end{align}
In the context of $S(q)$, they \cite[Remark 2]{AB} remarked that
\begin{quote}
	{\it It is therefore natural to expect modular or mock-modular behaviour of $S(q)$.}	
\end{quote}
 According to Ramanujan \cite{BB}, a {\it mock theta function}\footnote{In modern language, a mock theta function is the holomorphic part of a so-called harmonic Maass form (a certain non-holomorphic modular object).} is a function $F$ that has asymptotic expansions at every roots of unity of the same type as those of weakly holomorphic modular forms, however there is not a single weakly holomorphic modular form whose asymptotic expansions agree at all roots of unity with that of $F$. %In his thesis \cite{Zw}, Zwegers ``completed" Ramanujan's mock theta functions to harmonic Maass forms and its holomorphic part is called a {\it mock modular form}.

In this paper, we confirm this by relating $S(q)$ to mock theta functions. Recall Ramanujan's third order mock theta function $\o(q)$ %\cite{AB5}
\begin{equation}\label{Ramanujanomega}
\omega(q):=\sum_{n\ge0}\frac{q^{2n(n+1)}}{\left(q;q^2\right)^2_{n+1}}=:\sum_{n\ge 0}c_{\o}(n)q^n
\end{equation}
and  McIntosh's \cite{McIntosh} second order mock theta function $B(q)$ 
\begin{equation}\label{McIntosh}
B(q):=\sum_{n\ge 0}\frac{\left(-q^2;q^2\right)_n q^{n(n+1)}}{\left(q;q^2\right)^2_{n+1}}=:\sum_{n\ge 0}c_B(n)q^n.
\end{equation}
\begin{theorem}\label{lem1}
	We have
	\[
	S(q)=2B(-q)-\frac{\left(q;q^2\right)^2_{\infty}}{\left(-q^2;q^2\right)_{\infty}}\o(-q).
	\]
\end{theorem}
 Similar to the Ramanujan’s congruences for the partition function
\begin{align*}
p(5n+4)&\equiv 0\pmod{5},\\ 
p(7n+5)&\equiv 0\pmod{7},\\ 
p(11n+6)&\equiv 0\pmod{11}.
\end{align*}
 Andrews and the third author conjectured the following \cite[Conjectures 1 and 2]{AB}.
\begin{conjecture}\label{conj1a}
For $n\in \mathbb{N}_0$, we have
\begin{align*}
c(8n+4)\equiv 0\pmod{4}.
\end{align*}
\end{conjecture}
\begin{conjecture}\label{conj1b}
	For $n\in \mathbb{N}_0$, we have
	\begin{align*}
	c(8n+6)\equiv 0\pmod{8}.
	\end{align*}
\end{conjecture}


In this paper, we prove these conjectures.

\begin{theorem}\label{ABConj1}
\Cref{conj1a} is true.
\end{theorem}

\begin{theorem}\label{ABConj1a}
	\Cref{conj1b} is true.
\end{theorem}

 Moreover, we find and prove the following congruence for $c(n)$. 

\begin{theorem}\label{newthm}
For $n\in \mathbb{N}_0$, we have
\[
	c(16n+13)\equiv 0 \pmod 4.
\]
\end{theorem}


In context of $S(q)$, Andrews and the third author proposed \cite[Conjecture 3]{AB} the following conjecture: 
\begin{conjecture}\label{Conj2}
For $n\in \mathbb{N}_0$, we have 
\[s(4n+1)=0.\]
\end{conjecture}

In this paper, we confirm this conjecture.

\begin{theorem}\label{ABConj2}
	\Cref{Conj2} is true.
\end{theorem}

%\begin{conjecture}\label{ABConj3}
%	For $N\in \mathbb{N}_0$, we have 
%\begin{align*}
%\textnormal{coeff}_{\left[q^{2N}\right]}\left(\left(q^4;q^4\right)_{\infty}S(q)\right)&\\[-1 em]
%&-\hspace{-1.75 cm}=\textnormal{coeff}_{\left[q^{2N}\right]}\left(\sum_{n\ge 0}(-1)^nq^{6n^2+4n}\left(1+q^{4n+2}\right)\sum_{|j|\le n}(-1)^{-j}q^{-2j^2}\right).
%\end{align*}
%\end{conjecture}



%{\bf \color{blue} K: added the following section. Placed the references used therein.}


%{\bf \color{blue} K: added organization.}

The remainder of the paper is organized as follows. In \Cref{sec1} we recall some well-known $q$-series transformations and identities related to modular forms and mock theta functions. In \Cref{sec2} we relate $S(q)$ to mock theta functions. In \Cref{sec3}, \Cref{sec3a}, and \Cref{sec3b}, we prove \Cref{ABConj1}, \Cref{ABConj1a}, and \Cref{newthm}, respectively. In \Cref{sec4} we prove \Cref{ABConj2}. Finally in \Cref{sec5} we give few follow up questions for future research.

%\section*{Acknowledgements}
%The first and the second author have received funding from the European Research Council (ERC)
%under the European Union's Horizon 2020 research and innovation programme (grant agreement
%No. 101001179).


\section{Preliminaries}\label{sec1}


In this section, we recall $q$-series identities and congruences for theta functions and mock theta functions.

We start with the following identity due to Ramanujan \cite[equation (3.17)$_\text{R}$]{A0}.

%{\bf KB: Sometimes we list conditions and sometimes not. \color{blue} K: Yes, because, in the references whenever conditions are mentioned, we only state the conditions, otherwise we did not. \color{black} KB: I feel should be consistent though. \color{blue} K: according to our discussion, let it be as it is.}

\begin{lemma}\label{lemident1} 
	We have %From \cite[]{A0}, we have
	\begin{align*}
	&\sum_{n\ge 0}\!\frac{\left(-aq,-bq\right)_nq^{n+1}}{\left(-cq\right)_n}\!
	=\!\frac{c}{ab}\!\sum_{n\ge 1}\!\frac{\left(-c^{-1}\right)_n \pa{ab}{c}^n\! q^{\frac{n(n+1)}{2}}}{\left(\frac{aq}{c},\frac{bq}{c}\right)_n}\\
	&\hspace{7 cm}-\!\frac{c\left(-aq,-bq\right)_{\infty}}{ab\left(-cq\right)_{\infty}}\!\sum_{n\ge 1}\!\frac{\pa{ab}{c^2}^nq^{n^2}}{\left(\frac{aq}{c},\frac{bq}{c}\right)_n}.%\label{Ident1}\\
	%&\sum_{n\ge0}\frac{(cq)_nq^n}{(-aq,-bq)_n}\!=\!\left(1+a^{-1}\right)\!\sum_{n\ge0}\!\frac{(cq)_n\left(-\frac{b}{a}\right)^n\!q^{\frac{n(n+1)}{2}}}{\left(-\frac{cq}{a}\right)_{n+1}(-bq)_n}\!-\!\frac{a^{-1}(cq)_{\infty}}{(-aq,-bq)_{\infty}}\!\sum_{n\ge0}\!\frac{\left(-\frac{b}{a}\right)^n\!q^{\frac{n(n+1)}{2}}}{\left(-\frac{cq}{a}\right)_{n+1}}.\label{ident2}
	\end{align*}
\end{lemma}

Recall the {\it theta function} %(see \cite[Theorem 1.60]{Ono}) {\bf \color{blue} K: I kept the ref. because here we do not speak about the Jacobi triple product identity. But if you feel we don't need it at all, I will delete it the next version.}
\begin{equation}\label{theta}
\Theta(q):=\sum_{n\in \mathbb{Z}}q^{n^2}=\frac{\left(q^2;q^2\right)^5_{\infty}}{(q)^2_{\infty}\left(q^4;q^4\right)^2_{\infty}}.
\end{equation}
Note that
\begin{equation}\label{thetas}
\Theta(-q)=\sum_{n\in \mathbb{Z}}(-1)^nq^{n^2}=\frac{(q)^2_{\infty}}{\left(q^2;q^2\right)_{\infty}}.
\end{equation}
Moreover, we require the ``odd" theta function, defined by 
\begin{equation}\label{psi}
\psi(q):=\frac{\left(q^2;q^2\right)^2_{\infty}}{\left(q\right)_{\infty}}=\sum_{n\ge 0}q^{\frac{n(n+1)}{2}}.
\end{equation}
Next, we state a result of Ramanujan, which simply follows by splitting into odd and even exponents.
\begin{lemma}\label{thetadissection}
	We have
	\[
	\Theta(q)=\Theta\left(q^4\right)+2q\psi\!\left(q^8\right).
	\]
\end{lemma} 

Due to Andrews \cite[(4.3)]{A}, we have the following Lerch-type representation of\,\footnote{The factor in the denominator in the sum on the right-hand side of \cite[(4.3)]{A} is $(q;q^2)_{n+1}$ (after fixing the typographical error $(q)_{2n+1}$).} $B(q)$. %(using \cite[(1.3)]{A}): %we have the following:

\begin{lemma}\label{And}
	We have %{\bf \color{blue} K: proof in the details file.}
	\[
	B(q)=\frac{\left(-q^2;q^2\right)_{\infty}}{\left(q^2;q^2\right)_{\infty}}\sum_{n\in \mathbb{Z}}\frac{(-1)^n q^{2n(n+1)}}{1-q^{2n+1}}.
	\]
\end{lemma} 

From \cite[(4.1)]{Mao}, we obtain that the restriction of $B(q)$ to exponents divisible by $4$ is a simple modular form.
\begin{lemma}\label{Mao}
	We have
	\[
	\sum_{n\ge 0}c_B(4n)q^n=\frac{\left(q^{2};q^2\right)^{14}_{\infty}}{(q)^9_{\infty} \left(q^4;q^4\right)^4_{\infty}}.
	\]
\end{lemma}

From Chan and Mao \cite[(2.13)]{CM}, we can conclude the following restriction of $B$ to exponents congruent to $1$ modulo $4$.
\begin{lemma}\label{CM}
	We have 
	\[
	\sum_{n\ge 0}c_B(4n+1)q^n%=2\frac{\left(-q\right)^2_{\infty}\left(q^4;q^4\right)^4_{\infty}\left(q^2;q^2\right)^2_{\infty}}{\left(q\right)^5_{\infty}\left(-q^2;q^2\right)^4_{\infty}}
	=\frac{2\left(q^2;q^2\right)^8_{\infty}}{(q)^7_{\infty}}. %2\frac{\left(-q;q^2\right)^2_{\infty}\left(q^4;q^4\right)^4_{\infty}}{\left(q\right)^3_{\infty}\left(q,q^3;q^4\right)^2_{\infty}\left(-q^2;q^2\right)^2_{\infty}}.
	\]
\end{lemma}




 Next, we require an identity between $\o(q)$ and $f(q)$, where Ramanujan's third order mock theta function $f(q)$ is defined by
\begin{equation*}%\label{Ramanujanf}
f(q):=\sum_{n\ge0}\frac{q^{n^2}}{\left(-q\right)^2_{n}}.
\end{equation*}

Due to Watson \cite[p. 66]{W} (see also \cite[(4)]{APSY}), we have the following identity. 
\begin{lemma}\label{odissection}
We have %{\bf \color{blue} K: changed $F\to G$ as in intro, mock theta function defn. used $F$.}
\[
f\!\left(q^8\right)+2q\o(q)+2q^3\o\!\left(-q^4\right)=\frac{\Theta(q)\Theta^2\left(q^2\right)}{\left(q^4;q^4\right)^2_{\infty}}:=G(q).
\]
\end{lemma}
%{\bf \color{blue} K: added the sentence.} Next, we state the $4$-dissection of $F(q)$.

Due to Andrews, Passary, Sellers, and Yee \cite[Lemma 3.1]{APSY}, we have the following dissection of $G(q)$.
\begin{lemma}\label{Fdissection}
We have
\[
G\left(q\right)=\sum_{j=0}^{4} q^jG_j\!\left(q^4\right),
\]
where
\begin{align*}
G_0(q)&:=\frac{\Theta^3(q)}{(q)^2_{\infty}},\hspace{2.6 cm}  G_1(q):=\frac{2\Theta^2(q)\psi\left(q^2\right)}{(q)^2_{\infty}},\\ G_2(q)&:=\frac{4\Theta(q)\psi^2\!\left(q^2\right)}{(q)^2_{\infty}},\hspace{1.5 cm} G_3(q):=\frac{8\psi^3\!\left(q^2\right)}{(q)^2_{\infty}}.
\end{align*}
\end{lemma}

Finally, we state congruences for $c_{\o}(n)$ \cite[Theorem 3.4]{APSY}. 
\begin{lemma}\label{Ref1}
	For $n\in \mathbb{N}_0$, we have 
	\begin{align*}
	c_{\o}\left(8n+3\right)&\equiv 0\!\pmod{4},\\
	c_{\o}\left(8n+5\right)&\equiv 0\!\pmod{8},\\
	c_{\o}\left(16n+12\right)&\equiv 0\!\pmod{4}.
	\end{align*}
\end{lemma}
 














\section{Proof of \Cref{lem1}}\label{sec2}
 
 In this section, we prove \Cref{lem1}.



\begin{proof}[Proof of \Cref{lem1}]
By \Cref{lemident1} with $q\mapsto q^2$, we have
\begin{align*}
\sum_{n\ge 0}\frac{\left(-aq^2,-bq^2;q^2\right)_nq^{2n+2}}{\left(-cq^2;q^2\right)_n}
&=\frac{c}{ab}\sum_{n\ge 1}\!\frac{\left(-c^{-1};q^2\right)_n \pa{ab}{c}^n q^{n(n+1)}}{\left(\frac{aq^2}{c},\frac{bq^2}{c};q^2\right)_n}\\
&\hspace{1 cm}-\frac{c\left(-aq^2,-bq^2;q^2\right)_{\infty}}{ab\left(-cq^2;q^2\right)_{\infty}}\sum_{n\ge 1}\!\frac{\pa{ab}{c^2}^nq^{2n^2}}{\left(\frac{aq^2}{c},\frac{bq^2}{c};q^2\right)_n}.
\end{align*}
Plugging $a=b=-q^{-1}$ and $c=1$ into the above, we obtain 
\begin{align*}
\sum_{n\ge 0}\frac{\left(q;q^2\right)^2_nq^{2n+2}}{\left(-q^2;q^2\right)_n}
%&=q^2\sum_{n\ge 1}\!\frac{\left(-1;q^2\right)_n q^{-2n} q^{n(n+1)}}{\left(-q;q^2\right)^2_n}-\frac{\left(q;q^2\right)^2_{\infty}q^2}{\left(-q^2;q^2\right)_{\infty}}\sum_{n\ge 1}\!\frac{q^{2n^2-2n}}{\left(-q;q^2\right)^2_n}\\
&=2q^2\sum_{n\ge 1}\!\frac{\left(-q^2;q^2\right)_{n-1} q^{n(n-1)}}{\left(-q;q^2\right)^2_n}-\frac{\left(q;q^2\right)^2_{\infty}q^2}{\left(-q^2;q^2\right)_{\infty}}\sum_{n\ge 1}\!\frac{q^{2n(n-1)}}{\left(-q;q^2\right)^2_n}\\
%\intertext{\bf KB: But for $\o$, there is $n+1$. \color{blue} K: fixed the typo: $n\to n+1$.}
&=2q^2\sum_{n\ge 0}\!\frac{\left(-q^2;q^2\right)_{n} q^{n(n+1)}}{\left(-q;q^2\right)^2_{n+1}}-\frac{\left(q;q^2\right)^2_{\infty}q^2}{\left(-q^2;q^2\right)_{\infty}}\sum_{n\ge 0}\!\frac{q^{2n(n+1)}}{\left(-q;q^2\right)^2_{n+1}}\\
&=2q^2B(-q)-\frac{\left(q;q^2\right)^2_{\infty}q^2}{\left(-q^2;q^2\right)_{\infty}}\o(-q),
\end{align*}
using \eqref{McIntosh} and \eqref{Ramanujanomega} with $q\mapsto -q$. Multiplying both sides by $q^{-2}$ and using \eqref{S}, the theorem follows.\qedhere 
\end{proof}


%{\bf \color{blue} K: new section added below.}

\section{Proof of \Cref{ABConj1}}\label{sec3}

 In this section, we prove \Cref{ABConj1}.

\begin{proof}[Proof of \Cref{ABConj1}]
Using \eqref{C}, \eqref{A}, and \Cref{lem1}, we have 
\begin{align}\label{eqn1}
C(q)%&=\frac{q\left(-q^2;q^2\right)_{\infty}}{\left(q;q^2\right)^2_{\infty}}\left(2B(-q)-\frac{\left(q;q^2\right)^2_{\infty}}{\left(-q^2;q^2\right)_{\infty}}\o(-q)\right)\nonumber\\
&=\frac{2q\left(-q^2;q^2\right)_{\infty}}{\left(q;q^2\right)^2_{\infty}}B(-q)-q\o(-q).
\end{align}
By \Cref{Ref1}, we have, for $n\in \mathbb{N}_0$, 
\begin{equation}\label{ref1}
\textnormal{coeff}_{\left[q^{8n+3}\right]}\left(\o(q)\right)\equiv 0\pmod{4}.
\end{equation}
Now, for a $q$-series $H(q)=\sum_{n\ge 0}h(n)q^n$ and $n\in \mathbb{N}_0$, we define
\[
\operatorname{coeff}_{\left[q^n\right]}\left(H(q)\right):=h(n).
\]
Note that \eqref{ref1} is equivalent to \[\textnormal{coeff}_{\left[q^{8n+4}\right]}\left(q\o(-q)\right)\equiv 0\pmod{4}.\]
% Define {\bf \color{blue} K: rewritten the defn. according to our definition using Andrews' notation $A(q)$.}
%\begin{equation}\label{AR}
%A(q):=\frac{A(q)}{q}=\frac{\left(-q^2;q^2\right)_{\infty}}{\left(q;q^2\right)^2_{\infty}},
%\end{equation}
 Thus, to prove \Cref{ABConj1}, by \eqref{A} and \eqref{eqn1}, we need to show that, for $n\in \mathbb{N}_0$,
\begin{align*}
%&\textnormal{coeff}_{\left[q^{8n+4}\right]}\!\left(2qA(q)B(-q)\right)\!\os{!}\equiv\! 0\!\pmod{4}\Leftrightarrow
 \textnormal{coeff}_{\left[q^{8n+3}\right]}\!\left(A(q)B(-q)\right)\!\equiv\! 0\!\pmod{2}.
\end{align*}

 Now, we reduce $A(q)B(-q)$ modulo $2$. First, we obtain, using (1.2), 
 \begin{align*}%\label{Amod2}
 A(q)=\frac{\left(-q^2;q^2\right)_{\infty}}{\left(q;q^2\right)^2_{\infty}}%\!\equiv\! \frac{\left(q^2;q^2\right)_{\infty}\left(q^2;q^2\right)^2_{\infty}}{\left(q\right)^2_{\infty}}
 &\!\equiv\!\left(q^2;q^2\right)^2_{\infty}\!\equiv\! \left(q^4;q^4\right)_{\infty}\!\pmod{2}.
 \end{align*}
 The final $q$-series only has exponents divisible by $4$.
 
 
 Next, we consider $B(-q)$. By \Cref{And}, we have 
 \begin{equation*}%\label{s1}
 B(-q)\equiv B(q)\equiv\sum_{n\in \mathbb{Z}}\frac{(-1)^n q^{2n(n+1)}}{1-q^{2n+1}}\pmod{2},
 \end{equation*}
using the fact that $\frac{(-q^2;q^2)_{\infty}}{(q^2;q^2)_{\infty}}\equiv 1\pmod{2}$. We then rewrite 
\begin{align*}%\label{s2}
\sum_{n\in \mathbb{Z}}\frac{(-1)^n q^{2n(n+1)}}{1-q^{2n+1}}&\equiv \sum_{n\in \mathbb{Z}}\frac{ q^{2n(n+1)}}{1+q^{2n+1}}=\frac 12\sum_{n\in \mathbb{Z}}\left(\frac{q^{2n(n+1)}}{1+q^{2n+1}}+\frac{q^{2n(n+1)}}{1+q^{-2n-1}}\right)\nonumber\\
&=\frac 12\sum_{n\in \mathbb{Z}}q^{2n(n+1)}\pmod{2}.
\end{align*}
Again, note that the final $q$-series only has exponents divisible by $4$. Combining gives the claim. %Combining this, \eqref{s1}, and \eqref{s2}, we have only even exponents of $\frac{(-q^2;q^2)_{\infty}}{(q;q^2)^2_{\infty}}B(-q)$ by reducing modulo $2$. 
This concludes the proof of \Cref{ABConj1}.\qedhere 
\end{proof}

%{\bf \color{blue} K: made new section.}

\section{Proof of \Cref{ABConj1a}}\label{sec3a}

In this section, we prove \Cref{ABConj1a}. 

\begin{proof}[Proof of \Cref{ABConj1a}]
 %Next, we prove the second congruence of \Cref{ABConj1}. 
We split the proof into two parts corresponding to the two terms in \Cref{lem1} (after multiplying the pre-factor, see \eqref{C}). The contribution of the mock theta function is handled using known congruences, while the remaining term is treated via a $4$-dissection and cancellation argument.

 By \Cref{Ref1}, we have, for $n\in \mathbb{N}_0$, %{\bf \color{blue} K: moved steps in details file.}
\[\textnormal{coeff}_{\bigl[q^{8n+6}\bigr]}\left(q\o(-q)\right)\equiv 0\pmod{8}.\]
Thus, to prove \Cref{ABConj1a}, by \eqref{eqn1}, \eqref{A}, and \Cref{Ref1}, we need to show that, for $n\in \mathbb{N}_0$, %{\bf \color{blue} K: first step in details file.}
\begin{align*}%\label{TS1}
%&\textnormal{coeff}_{\bigl[q^{8n+6}\bigr]}\!\left(2qA(q)B(-q)\right)\!\os{!}\equiv\! 0\!\pmod{8}\nonumber\\
%&\hspace{5 cm}\Leftrightarrow
 \textnormal{coeff}_{\bigl[q^{8n+5}\bigr]}\!\left(A(q)B(-q)\right)\!\equiv\! 0\!\pmod{4}.
\end{align*}
%Now we reduce
%\[
%A(q)B(-q)\pmod{4}.
%\] 
In the following we even prove the stronger result that %To prove \eqref{TS1}, It suffices to show that
\begin{equation}\label{TSalt}
\textnormal{coeff}_{\left[q^{4n+1}\right]}\left(A(q)B(-q)\right)\equiv 0\pmod{4}.
\end{equation}
 %This is because only the residue class $4n+1$ contributes to the desired congruence.
%{\bf \color{blue} K: deleted sieving operator.}% Recall the {\it sieving operator} $V_{N,r}$ (for $N\in \mathbb{N}, r\in \mathbb{Z}$) defined on a $q$-series $V(q)=\sum_{n\ge 0} v(n)q^n$ as 
%\[
%V\!\mid\!S_{N,r}(q):=\sum_{n\equiv r\pmod N}v(n) q^n.
%\]

We first show that\footnote{Andrews and Bachraoui \cite[Remark 3]{AB} noticed that \eqref{claim1} holds. However, we give a proof for the reader's convenience.} 
 \begin{equation}\label{claim1}
\textnormal{coeff}_{\left[q^{4n+j}\right]}\left(A(q)\right)\equiv 0\pmod{4}\ \ \text{for}\ \ j\in\{2,3\}.
 \end{equation}
For this definitions \eqref{A} and \eqref{thetas}, we may write 
\begin{equation}\label{claim1eqn1}
A(q)%=\frac{\left(-q^2;q^2\right)_{\infty}\left(q^2;q^2\right)^2_{\infty}}{\left(q\right)^2_{\infty}}=\frac{\left(q^4;q^4\right)_{\infty}\left(q^2;q^2\right)_{\infty}}{\left(q\right)^2_{\infty}}\os{\eqref{thetas}}
=\frac{\left(q^4;q^4\right)_{\infty}}{\Theta(-q)}.
\end{equation}
Now, by \eqref{thetas}, note that 
\[
\Theta(-q)\equiv 1+2T(q)\pmod{4},
\]
where $T(q):=\sum_{n\ge 1} q^{n^2}$. Plugging this into \eqref{claim1eqn1}, we obtain 
\begin{align}\label{claim1eqn2}
\hspace{-0.33 cm}A(q)&\equiv\!\frac{\left(q^4;q^4\right)_{\infty}}{1+2T(q)}\!\equiv\! \left(q^4;q^4\right)_{\infty}\!\left(1-2T(q)\right)\!\equiv\! \left(q^4;q^4\right)_{\infty}\!\left(1+2T(q)\right)\!\!\pmod{4}.
\end{align}
Next, we write %using definitions \eqref{theta} and \eqref{psi}, note that \Cref{thetadissection} implies 
\[
T(q)=%\sum_{n\ge 1}q^{4n^2}+\sum_{n\ge 0}q^{(2n+1)^2}=
\sum_{n\ge 1}q^{4n^2}+q\sum_{n\ge 0}q^{4n(n+1)}.
\]
Plugging this into \eqref{claim1eqn2}, we obtain
\begin{equation}\label{claim1eqn3}
A(q)\equiv \left(q^4;q^4\right)_{\infty}\sum_{n\in \mathbb{Z}}q^{4n^2}+2q\left(q^4;q^4\right)_{\infty}\sum_{n\ge 0}q^{4n(n+1)}\pmod{4}.
\end{equation}
This proves \eqref{claim1}. 

Next, by \eqref{claim1eqn3}, we have 
\begin{equation*}%\label{Alevel}
A(q)\equiv A_0\!\left(q^4\right)+qA_1\!\left(q^4\right)\pmod{4},
\end{equation*}
where 
\begin{align}\label{claim2eqn1}
A_0(q)&:=\left(q\right)_{\infty}\sum_{n\in \mathbb{Z}}q^{n^2}=\frac{\left(q^2;q^2\right)^5_{\infty}}{\left(q\right)_{\infty}\left(q^{4};q^{4}\right)^2_{\infty}},\nonumber\\
A_1(q)&:=2\left(q\right)_{\infty}\sum_{n\ge 0}q^{n(n+1)}=\frac{2\left(q\right)_{\infty}\left(q^{4};q^{4}\right)^2_{\infty}}{\left(q^2;q^2\right)_{\infty}},
\end{align}
using \eqref{theta} and \eqref{psi}. 

Now write
\[
B(q)=\sum_{j=0}^{3}q^j\!B_j\!\left(q^4\right).
\]
Then %{\bf \color{blue} K: corrected: $B(q)\to B(-q)$.}
\[
A(q)B(-q)\equiv \left(A_0\!\left(q^4\right)+qA_1\!\left(q^4\right)\right)\sum_{j=0}^{3}(-1)^jq^j\!B_j\!\left(q^4\right)\!\pmod{4}=:\sum_{j=0}^3q^j\!\mathcal{F}_j\!\left(q^4\right),
\]
for certain $q$-series $\mathcal{F}_j$. In particular, %{\bf \color{blue} K: placed $-$ sign there.} 
\[
\mathcal{F}_1(q)=-A_0\!\left(q\right)B_1\!\left(q\right)+A_1\!\left(q\right)B_0\!\left(q\right).
\]
To finish the proof of \eqref{TSalt}, we need to show that
\begin{equation}\label{F1}
\mathcal{F}_1(q)\equiv 0\pmod{4}.
\end{equation}
We now determine $B_0$ and $B_1$. For this, using \Cref{Mao} and \Cref{CM}, we note % Next, note that to prove \eqref{TSalt}, we only need $0,1\pmod{4}$ residue classes for $B$. We isolate residue classes $\pmod{4}$ using the sieve operator $S_{4,j}$, and analyze each component separately. Precisely, we have 
%\begin{align*}
%B(q)\mid S_{4,1}&=\sum_{n\ge 0}\b(4n+1)q^{4n+1}\os{\substack{\text{\Cref{CM}}\\(q\mapsto q^4)}}=2q\frac{\left(q^8;q^8\right)^8_{\infty}}{\left(q^4;q^4\right)^7_{\infty}},\\%2q\frac{\left(-q^4;q^8\right)^2_{\infty}\left(q^{16};q^{16}\right)^4_{\infty}}{\left(q^4;q^4\right)^3_{\infty}\left(q^4,q^{12};q^{16}\right)^2_{\infty}\left(-q^8;q^8\right)^2_{\infty}},\\
%B(q)\mid S_{4,0}&=\sum_{n\ge 0}\b(4n)q^{4n}\os{\substack{\text{\Cref{Mao}}\\(q\mapsto q^4)}}=\frac{\left(q^{8};q^8\right)^{14}_{\infty}}{\left(q^4;q^4\right)^9_{\infty} \left(q^{16};q^{16}\right)^4_{\infty}}.
%\end{align*}
%Thus 
\begin{align}\label{claim2eqn2}
B_0(q)&=\sum_{n\ge 0}c_B(4n)q^n=\frac{\left(q^{2};q^2\right)^{14}_{\infty}}{\left(q\right)^9_{\infty} \left(q^{4};q^{4}\right)^4_{\infty}},\nonumber\\%\frac{\left(-q^4;q^8\right)^2_{\infty}\left(q^{16};q^{16}\right)^4_{\infty}}{\left(q^4;q^4\right)^3_{\infty}\left(q^4,q^{12};q^{16}\right)^2_{\infty}\left(-q^8;q^8\right)^2_{\infty}},\nonumber\\
B_1(q)&=\sum_{n\ge 0}c_B(4n+1)q^n=-\frac{2\left(q^2;q^2\right)^8_{\infty}}{\left(q\right)^7_{\infty}}.
\end{align}
 Combining \eqref{claim2eqn1} and \eqref{claim2eqn2}, we obtain 
 \begin{align*}
 %\mathcal{A}^{[1]}_{4}(q)B^{[0]}_4(q)&=2q\left(q^4;q^4\right)_{\infty}\frac{\left(q^{16};q^{16}\right)^2_{\infty}}{\left(q^8;q^8\right)_{\infty}}\frac{\left(q^{8};q^8\right)^{14}_{\infty}}{\left(q^4;q^4\right)^9_{\infty} \left(q^{16};q^{16}\right)^4_{\infty}}\\
 %&=2q\frac{\left(q^{8};q^8\right)^{13}_{\infty}}{\left(q^4;q^4\right)^8_{\infty} \left(q^{16};q^{16}\right)^2_{\infty}},\\
 %\intertext{\bf \color{blue} K: rewritten the following.}
 A_0(q)B_1(q)%&=-2\frac{\left(q^2;q^2\right)^5_{\infty}\left(q^2;q^2\right)^8_{\infty}}{\left(q\right)_{\infty}\left(q^{4};q^{4}\right)^2_{\infty}\left(q\right)^7_{\infty}}
 =-\frac{2\left(q^{2};q^2\right)^{13}_{\infty}}{\left(q\right)^8_{\infty} \left(q^{4};q^{4}\right)^2_{\infty}}.%\frac{\left(q^8;q^8\right)^5_{\infty}}{\left(q^4;q^4\right)_{\infty}\left(q^{16};q^{16}\right)^2_{\infty}}\frac{\left(-q^4;q^8\right)^2_{\infty}\left(q^{16};q^{16}\right)^4_{\infty}}{\left(q^4;q^4\right)^3_{\infty}\left(q^4,q^{12};q^{16}\right)^2_{\infty}\left(-q^8;q^8\right)^2_{\infty}}\\
 %&=-2q \frac{\left(-q^4;q^8\right)^2_{\infty}\left(q^8;q^8\right)^5_{\infty}\left(q^{16};q^{16}\right)^2_{\infty}}{\left(q^4;q^4\right)^4_{\infty}\left(q^4,q^{12};q^{16}\right)^2_{\infty}\left(-q^8;q^8\right)^2_{\infty}}\\
 %&=-2q\frac{\left(-q^4;q^8\right)^2_{\infty}\left(q^8;q^8\right)^5_{\infty}\left(q^{16};q^{16}\right)^4_{\infty}\left(q^8;q^{16}\right)^2}{\left(q^4;q^4\right)^6_{\infty}\left(-q^8;q^8\right)^2_{\infty}}\\
 %&=-2q\frac{\left(-q^4;q^8\right)^2_{\infty}\left(q^8;q^8\right)^7_{\infty}\left(q^{16};q^{16}\right)^2_{\infty}}{\left(q^4;q^4\right)^6_{\infty}\left(-q^8;q^8\right)^2_{\infty}}=-2q\frac{\left(-q^4;q^8\right)^2_{\infty}\left(q^8;q^8\right)^9_{\infty}}{\left(q^4;q^4\right)^6_{\infty}}\\
 %&=-2q\frac{\left(q^8;q^{16}\right)^2_{\infty}\left(q^8;q^8\right)^9_{\infty}}{\left(q^4;q^8\right)^2_{\infty}\left(q^4;q^4\right)^6_{\infty}}=-2q\frac{\left(q^8;q^8\right)^{11}_{\infty}}{\left(q^4;q^8\right)^2_{\infty}\left(q^4;q^4\right)^6_{\infty}\left(q^{16};q^{16}\right)^2_{\infty}}\\
 %&=-2q\frac{\left(q^8;q^8\right)^{13}_{\infty}}{\left(q^4;q^4\right)^8_{\infty}\left(q^{16};q^{16}\right)^2_{\infty}}.
 \end{align*}
Similarly, combining \eqref{claim2eqn1} and \eqref{claim2eqn2}, we obtain 
\begin{align*}
A_1(q)B_0(q)%&=2\left(q\right)_{\infty}\frac{\left(q^{4};q^{4}\right)^2_{\infty}}{\left(q^2;q^2\right)_{\infty}}\frac{\left(q^{2};q^2\right)^{14}_{\infty}}{\left(q\right)^9_{\infty} \left(q^{4};q^{4}\right)^4_{\infty}}%\\
=\frac{2\left(q^{2};q^2\right)^{13}_{\infty}}{\left(q\right)^8_{\infty} \left(q^{4};q^{4}\right)^2_{\infty}}.\\
%\intertext{\bf \color{blue} K: rewritten the following.}
%\mathcal{A}^{[0]}_{4}(q)B^{[1]}_4(q)&=-2q\frac{\left(q^8;q^8\right)^5_{\infty}\left(q^8;q^8\right)^8_{\infty}}{\left(q^4;q^4\right)_{\infty}\left(q^{16};q^{16}\right)^2_{\infty}\left(q^4;q^4\right)^7_{\infty}}=-2q\frac{\left(q^{8};q^8\right)^{13}_{\infty}}{\left(q^4;q^4\right)^8_{\infty} \left(q^{16};q^{16}\right)^2_{\infty}}.%\frac{\left(q^8;q^8\right)^5_{\infty}}{\left(q^4;q^4\right)_{\infty}\left(q^{16};q^{16}\right)^2_{\infty}}\frac{\left(-q^4;q^8\right)^2_{\infty}\left(q^{16};q^{16}\right)^4_{\infty}}{\left(q^4;q^4\right)^3_{\infty}\left(q^4,q^{12};q^{16}\right)^2_{\infty}\left(-q^8;q^8\right)^2_{\infty}}\\
%&=-2q \frac{\left(-q^4;q^8\right)^2_{\infty}\left(q^8;q^8\right)^5_{\infty}\left(q^{16};q^{16}\right)^2_{\infty}}{\left(q^4;q^4\right)^4_{\infty}\left(q^4,q^{12};q^{16}\right)^2_{\infty}\left(-q^8;q^8\right)^2_{\infty}}\\
%&=-2q\frac{\left(-q^4;q^8\right)^2_{\infty}\left(q^8;q^8\right)^5_{\infty}\left(q^{16};q^{16}\right)^4_{\infty}\left(q^8;q^{16}\right)^2}{\left(q^4;q^4\right)^6_{\infty}\left(-q^8;q^8\right)^2_{\infty}}\\
%&=-2q\frac{\left(-q^4;q^8\right)^2_{\infty}\left(q^8;q^8\right)^7_{\infty}\left(q^{16};q^{16}\right)^2_{\infty}}{\left(q^4;q^4\right)^6_{\infty}\left(-q^8;q^8\right)^2_{\infty}}=-2q\frac{\left(-q^4;q^8\right)^2_{\infty}\left(q^8;q^8\right)^9_{\infty}}{\left(q^4;q^4\right)^6_{\infty}}\\
%&=-2q\frac{\left(q^8;q^{16}\right)^2_{\infty}\left(q^8;q^8\right)^9_{\infty}}{\left(q^4;q^8\right)^2_{\infty}\left(q^4;q^4\right)^6_{\infty}}=-2q\frac{\left(q^8;q^8\right)^{11}_{\infty}}{\left(q^4;q^8\right)^2_{\infty}\left(q^4;q^4\right)^6_{\infty}\left(q^{16};q^{16}\right)^2_{\infty}}\\
%&=-2q\frac{\left(q^8;q^8\right)^{13}_{\infty}}{\left(q^4;q^4\right)^8_{\infty}\left(q^{16};q^{16}\right)^2_{\infty}}.
\end{align*}
Thus \eqref{F1} holds.\qedhere 
%Therefore,
%\begin{equation}\label{final}
%\mathcal{A}^{[1]}_{4}(q)B^{[0]}_4(q)+\mathcal{A}^{[0]}_{4}(q)B^{[1]}_4(q)=0.
%\end{equation}
% The crucial point is that the mixed terms cancel $\pmod{4}$, which eliminates all contributions in the residue class $4n+1$. Using \eqref{Alevel}, we note that, to prove \eqref{TSalt}, it suffices to show that
%\begin{align*}
%&A(q)\mid S_{4,0}\,B(-q)\mid S_{4,1}+A(q)\mid S_{4,1}\,B(-q)\mid S_{4,0}\\
%&\equiv \underset{\os{\eqref{claim2eqn2}}=\mathcal{A}^{[0]}_{4}(q)B^{[1]}_4(q)+\mathcal{A}^{[1]}_{4}(q)B^{[0]}_4(q)}{\underbrace{\mathcal{A}^{[0]}_{4}(q)\left(B(-q)\right)\mid S_{4,1}+\mathcal{A}^{[1]}_{4}(q) \left(B(-q)\right)\mid S_{4,0}}}\pmod{4}\os{!}=0.
%\end{align*}
%This follows from the cancellation in \eqref{final}.\qedhere 
\end{proof}

%{\bf \color{blue} K: New section added below to include the proof of \Cref{newthm}.}

\section{Proof of \Cref{newthm}}\label{sec3b}

In this section, we prove \Cref{newthm}. First, we show the following.

\begin{proposition}\label{newprop}
	We have
	\[
	A(q)B(-q)\equiv (q^{16};q^{16})_\infty \pmod 2.
	\]
\end{proposition}

\begin{proof}
	First, by \eqref{A}, we have
	\[
	A(q)=\frac{(-q^2;q^2)_\infty}{(q;q^2)_\infty^2}
	\equiv \frac{(q^2;q^2)_\infty}{(q^2;q^4)_\infty}
	=(q^4;q^4)_\infty \pmod 2.
	\]
	Next, \Cref{And} gives
	\[
	B(-q)=\frac{(-q^2;q^2)_\infty}{(q^2;q^2)_\infty}
	\sum_{n\in \mathbb{Z}}\frac{(-1)^n q^{2n(n+1)}}{1+q^{2n+1}}
	\equiv
	\sum_{n\in \mathbb{Z}}\frac{q^{2n(n+1)}}{1+q^{2n+1}} \pmod 2.
	\]
	Pairing the terms for $n$ and $-n-1$, we obtain
	\[
	\frac{q^{2n(n+1)}}{1+q^{2n+1}}
	+
	\frac{q^{2n(n+1)}}{1+q^{-2n-1}}
	= q^{2n(n+1)}.
	\]
	Hence
	\[
	B(-q)\equiv \sum_{n\ge 0} q^{2n(n+1)} = \psi(q^4) \pmod 2.
	\]
	Using \eqref{psi}, we conclude that
	\[
	A(q)B(-q)
	\equiv \left(q^4;q^4\right)_\infty\,\psi\!\left(q^4\right)
	= \left(q^8;q^8\right)_\infty^2
	\equiv \left(q^{16};q^{16}\right)_\infty \pmod 2.\qedhere
	\]
%	since $(1-q^{8m})^2\equiv 1-q^{16m} \pmod 2$ for $m\in \mathbb{N}$.
\end{proof}

Now we are ready to prove \Cref{newthm}.

\begin{proof}[Proof of \Cref{newthm}]
By \Cref{newprop}, we have
\begin{equation}\label{neqn1}
2qA(q)B(-q)\equiv 2q\left(q^{16};q^{16}\right)_\infty \pmod 4.
\end{equation}
 Since $\textnormal{coeff}_{[q^{16n+j}]}(q(q^{16};q^{16})_\infty)=0$ for $j\neq 1$, by \eqref{neqn1},
$\textnormal{coeff}_{[q^{16n+13}]}(2qA(q)$\ $B(-q))\equiv 0\pmod{4}$. On the other hand, 
$\textnormal{coeff}_{[q^{16n+13}]}(q\omega(-q))=c_{\omega}(16n+12)$, because $(-1)^{16n+12}=1$.
By \Cref{Ref1}, we obtain
\[
c_{\omega}(16n+12)\equiv 0 \pmod 4.
\]
Therefore, using \eqref{eqn1}, we conclude the claim.\qedhere
%\[
%c(16n+13)\equiv 0 \pmod 4.\qedhere 
%\]	
\end{proof}

%{\bf \color{blue} K: up to here, added the calculations of Mohamed.}


%Define
%\begin{equation}\label{AR}
%A(q):=\frac{\left(-q^2;q^2\right)_{\infty}}{\left(q;q^2\right)^2_{\infty}},\quad \quad R(q):=A(q)B(-q).
%\end{equation}
%Note that
%\begin{equation*}%\label{ref1}
%A(-q)=\frac{\left(-q^2;q^2\right)_{\infty}}{\left(-q;q^2\right)^2_{\infty}}\equiv \frac{\left(-q^2;q^2\right)_{\infty}}{\left(q;q^2\right)^2_{\infty}}=A(q)\pmod 4,
%\end{equation*}
%using $(1-x)^2\equiv (1+x)^2\pmod{4}$ in the second step. Thus, we have
%\begin{equation}\label{ref1}
%A(q)\equiv A(-q)\pmod{4}.
%\end{equation} 
%We can write 
%\begin{equation}\label{A1}
%A(q)=A_0\!\left(q^2\right)+2qA_{1}\!\left(q^2\right).
%\end{equation}
%Indeed, we have
%\[
%A(q)=A_0\!\left(q^2\right)+qA^*_1\left(q^2\right).
%\]
%We have to show that $2\mid A^*_1(q)$. Using \eqref{ref1}, we have
%\begin{align*}
%\hspace{-5 cm}A(-q)&=A_0\!\left(q^2\right)-qA^*_1\!\left(q^2\right)\os{!}\equiv A_0\!\left(q^2\right)+qA^*_1\!\left(q^2\right)\pmod{4}\\
%&\Leftrightarrow -A^*_1\!\left(q^2\right)\os{!}\equiv A^*_1\!\left(q^2\right)\pmod{4} \Leftrightarrow 2A^*_1\!\left(q^2\right)\os{!}\equiv 0\pmod{4}  \Leftrightarrow A^*_1\!\left(q^2\right)\equiv 0\pmod{2}. 
%\end{align*}
%We now apply the same argument for $B$. Precisely, by \eqref{McIntosh}, we have
%\[
%B(-q)=\sum_{n\ge 0}\frac{\left(-q^2;q^2\right)_n (-q)^{n(n+1)}}{\left(-q;q^2\right)^2_{n+1}}\equiv \sum_{n\ge 0}\frac{\left(-q^2;q^2\right)_n q^{n(n+1)}}{\left(q;q^2\right)^2_{n+1}}=B(q)
%\]
%using $(1-x)^2\equiv (1+x)^2\pmod{4}$. Thus, we have
%\[
%B(-q)\equiv B(q)\pmod{4}.
%\]
%Consequently, we can write 
%\begin{equation}\label{B1}
%B(q)=B_0\!\left(q^2\right)+2qB_1\!\left(q^2\right).
%\end{equation}

%Thus, we obtain
%\begin{align}\label{RDiss}
%R(q)&\os{\eqref{AR}}=A(q)B(-q)\os{\substack{\eqref{A1}\\\eqref{B1}}}=\left(A_0\!\left(q^2\right)+2qA_1\!\left(q^2\right)\right)\left(B_0\!\left(q^2\right)-2qB_1\!\left(q^2\right)\right)\nonumber\\
%&\hspace{0.1 cm}\equiv \left(A_0\!\left(q^2\right)+2qA_1\!\left(q^2\right)\right)\left(B_0\!\left(q^2\right)+2qB_1\!\left(q^2\right)\right)\nonumber\\
%&\hspace{0.1 cm}\equiv\underset{\text{even exponent}}{\underbrace{A_0\!\left(q^2\right)\!B_0\!\left(q^2\right)}}+2q\Bigl(A_0\!\left(q^2\right)\!B_1\!\left(q^2\right)+A_1\!\left(q^2\right)\!B_0\!\left(q^2\right)\Bigr)\pmod{4}.
%\end{align}
%To prove \eqref{TS1}, we need to show
%\begin{equation}\label{TS2}
%\textnormal{coeff}_{\bigl[q^{8n+4}\bigr]}\left(A_0\!\left(q^2\right)\!B_1\!\left(q^2\right)+A_1\!\left(q^2\right)\!B_0\!\left(q^2\right)\right)\os{!}\equiv 0\pmod{2}.
%\end{equation}

%Next, we note that
%\[
%A(q)\os{\eqref{A1}}\equiv A_0\!\left(q^2\right)\pmod{2}\quad \quad \text{and}\quad \quad A(q)\os{\eqref{Amod2}}\equiv \left(q^4;q^4\right)_{\infty}\pmod{2}.
%\]
%Thus
%\begin{equation}\label{A0}
%A_0\!\left(q^2\right)\equiv \left(q^4;q^4\right)_{\infty}\pmod{2}.










%Define
%\[
%E(k,m,n):=k(3k-1)+n(n+1)+2nm+n+m\quad \text{and}\quad X_{k,m}:=\left\{n\in \mathbb{N}_0:E(k,m,n)\equiv 2\pmod{4}\right\}.
%\]
%\begin{proposition}\label{prop0}
%Fix $k\in \mathbb{Z}$ and $m\in \mathbb{N}_0$. Then there exists a fixed point free involution map $\Phi: X_{k,m}\to X_{k,m}$. 	
%\end{proposition}
%\begin{proof}
%Given $k\in \mathbb{Z}$ and $m\in \mathbb{N}_0$, we know that (shown above)
%\[
%n\in S_{k,m}\Leftrightarrow n\equiv m\pmod{2}.
%\]
%Setting $m:=2s+\varepsilon$ with $\varepsilon\in \{0,1\}$, we have $n=2r+\varepsilon$. So,
%\begin{align}\label{6}
%E(k,2s+\varepsilon,2r+\varepsilon)
%&=k(3k-1)+(2r+\varepsilon)^2+2(2r+\varepsilon)
%+2(2r+\varepsilon)(2s+\varepsilon)+2s+\varepsilon \nonumber\\
%&=k(3k-1)+4r^2+4(2\varepsilon+1)r+8rs+2(2\varepsilon+1)s+3\varepsilon(\varepsilon+1)\nonumber\\
%&= k(3k-1)+4r\bigl(r+2s+2\varepsilon+1\bigr)
%+2(2\varepsilon+1)s+3\varepsilon(\varepsilon+1).
%\end{align}
%\end{proof}

%\begin{proposition}\label{prop1}
%If for any fixed $k\in \mathbb{Z}$ and $m\in \mathbb{N}_0$, there exists an $n\in \mathbb{N}_0$ such that $E(k,n,m)\equiv 2\pmod{4}$ and a fixed point free involution map $\Phi:X_{k,m}\to X_{k,m}$, then 
%\[
%\textnormal{coeff}_{\bigl[q^{4N+2}\bigr]}\left(\sum_{\substack{k\in \mathbb{Z}\\m\ge 0\\n\ge 0}}q^{E(k,n,m)}\right)\equiv 0\pmod{2}.
%\]
%\end{proposition}
%\begin{proof}
% Note that to prove the claim, we need to show that
%\[
%\sum_{\substack{k\in \mathbb{Z}\\m\ge 0\\n\ge 0}}q^{E(k,m,n)}\mid S_{4,2}\os{!}\equiv 0\pmod{2}\Leftrightarrow \sum_{\substack{k\in \mathbb{Z}\\m\ge 0}}\sum_{n\in X_{k,m}}q^{E(k,m,n)}\os{!}\equiv 0\pmod{2}.
%\]
%Since $\Phi$ is a fixed point free involution map on $X_{k,m}$, thus, for each $n\in X_{k,m}$, we have unique $\Phi(n)$ with $E(k,m,\Phi(n))\equiv 2\pmod{4}$. Consequently, for fixed $k\in \mathbb{Z},m\in \mathbb{N}_0$, we obtain
%\[
%\sum_{n\in X_{k,m}}q^{E(k,m,n)}\equiv 0\pmod{2}.
%\] 
%Therefore, summing over $k\in \mathbb{Z}$ and $m\in \mathbb{N}_0$, we conclude the proposition.
%\end{proof}


%To prove this, we show for $N\in \mathbb{N}_0$ fixed, 
%\begin{align}\label{setcheck}
%&\left|\left\{2k(3k-1)+2n(n+1)+(2n+1)(2m+1)-1=8N+4:k\in \mathbb{Z}, m,n\in \mathbb{N}_0\right\}\right|\os{!}\equiv\! 0\!\pmod{2}\nonumber\\
%&\Leftrightarrow \left|\left\{n^2+2mn+3k^2+2n-k=4N+2:k\in \mathbb{Z}, m,n\in \mathbb{N}_0\right\}\right|\os{!}\equiv\! 0\!\pmod{2}.
%\end{align}
%Note that for $n$ odd, 
%\begin{equation}\label{oddcase}
%n^2+2mn+3k^2+2n-k\os{!}\equiv 2\pmod{4}\Leftrightarrow 2m-k(k+1)\os{!}\equiv 3\pmod{4},
%\end{equation}
%because $n^2\equiv 1\pmod{4}$, $2mn\equiv 2m\pmod{4}$, $2n\equiv 2\pmod{4}$, and $3k^2-k\equiv -k(k+1)\pmod{4}$. But since both $2m$ and $k(k+1)$ even, \eqref{oddcase} doesn't hold. Therefore, to prove \eqref{setcheck}, we only need to consider $n$ even and thus to show
%\begin{align}\label{setcheck2}
%&\left|\left\{n^2+2mn+3k^2+2n-k=4N+2:k\in \mathbb{Z}, m,n\in \mathbb{N}_0\ \text{with}\ n\equiv 0\pmod{2}\right\}\right|\os{!}\equiv\! 0\!\pmod{2}\nonumber\\
%&\Leftrightarrow\left|\left\{4n^2+4mn+3k^2+4n-k=4N+2:k\in \mathbb{Z}, m,n\in \mathbb{N}_0\right\}\right|\os{!}\equiv\! 0\!\pmod{2}.
%\end{align}
%Since $4\mid 4n^2, 4mn, 4n$, to prove \eqref{setcheck2}, we need to show
%\begin{align*}
%&3k^2-k\os{!}\equiv 2\pmod{4}\Leftrightarrow -k(k+1)\os{!}\equiv 2\pmod{4}\Leftrightarrow\frac 14\left((2k+1)^2-1\right)\os{!}\equiv 2\pmod{4}\\
%&\Leftrightarrow (2k+1)^2\os{!}\equiv 9\pmod{16}\Leftrightarrow 2k+1\os{!}\equiv \pm 3\pmod{4}\Leftrightarrow k\os{!}\equiv 1\pmod{4}\ \text{or}\ k\os{!}\equiv 2\pmod{4}.
%\end{align*} 
%{\bf \color{blue} K: Up to here.}
%To prove this, it suffices to show that for $n\in \mathbb{N}_0$ {\bf KB: I believe in $(q^2;q^2)_{\infty}$ the minus stays. \color{blue} K: Yes, corrected it.}
%\begin{equation}\label{eqv1}
%\textnormal{coeff}_{\left[q^{4n+3}\right]}\left(\frac{\left(-q^2;q^2\right)_{\infty}}{\left(-q;q^2\right)^2_{\infty}}B(q)\right)\os{!}\equiv 0\pmod{2}.
%\end{equation}

%{\bf KB: Not sure what you do here. \color{blue} K: I just define it for the proof. If not needed, I will delete it later.} For a power series $F(q):=\sum_{n\ge 0}a(n)q^n$, for $M\in \mathbb{N}_{\ge 2}$ and $0\le j\le M-1$, define $F^{[M]}_j(q):=a(Mn+j)q^{Mn_+j}$. Thus, $F=\sum_{j=0}^{M-1}F^{[M]}_j$. 

%Define $A(q):=\frac{(-q^2;q^2)_{\infty}}{(-q;q^2)^2_{\infty}}$ and write it (choosing $M=4$) as% $A(q)=\sum_{j=0}^{3}A^{[4]}_j(q)$ and $B(q)=\sum_{\ell=0}^{3}B^{[4]}_{\ell}(q)$. Therefore, \eqref{eqv1} is equivalent to 
%\begin{equation}\label{eqv2}
%\sum_{j=0}^{3}A^{[4]}_j(q)B^{[4]}_{3-j}(q)\os{!}\equiv 0\pmod{2}.
%\end{equation}
%Due to Chan and Mao \cite[Corollary 2.3]{CM}, we have $B^{[4]}_{\ell}\equiv 0\pmod{2}$ for $\ell\in\{1,2\}$. Thus to prove \eqref{eqv2}, we need show that
%\begin{equation*}
%\sum_{j\in\{0,3\}}A^{[4]}_j(q)B^{[4]}_{3-j}(q)\os{!}\equiv 0\pmod{2}.
%\end{equation*} 
%To prove this, it is enough to show that $A^{[4]}_3(q)\equiv 0\pmod{2}$ and $B^{[4]}_3(q)\equiv 0\pmod{2}$. Note that
%\[
%A^{[2]}_1(q)\equiv 0\pmod{2} \Rightarrow A^{[4]}_3(q)\equiv 0\pmod{2}\quad \text{and}\quad  B^{[2]}_1(q)\equiv 0\pmod{2} \Rightarrow B^{[4]}_3(q)\equiv 0\pmod{2}.
%\]
%Consequently, to conclude the proof, we need to show that
%\begin{equation}\label{ts1}
%A^{[2]}_1(q)\equiv 0\pmod{2}\quad \quad \text{and}\quad \quad B^{[2]}_1(q)\equiv 0\pmod{2}.
%\end{equation}
%We have
%\begin{align*}
%\frac{(-q^2;q^2)_{\infty}}{(-q;q^2)^2_{\infty}}\equiv \frac{(q^2;q^2)_{\infty}}{\left(q^2;q^4\right)_{\infty}}\pmod{2}=\left(q^4;q^4\right)\pmod{2}
%\end{align*}
%{\bf \color{blue} K: proof to be continued....}
%\end{proof}


\section{Proof of \Cref{ABConj2}}\label{sec4}

In this section, we prove \Cref{ABConj2}. Write 
\[
\Theta(q)\o(q)=:\sum_{n\ge 0}\d(n)q^n.
\]
 For a $q$-series $H(q):=\sum_{n\ge 0}h(n) q^n$, we define the operator 
\[
\mathcal{R}\left(H(q)\right):=\sum_{n\ge 0}h\!\left(4n+1\right) q^n.
\] and set 
%\[M(q):=\sum_{n\ge 0}\d\left(4n+1\right)q^n.\] 
\[
M(q):=\mathcal{R}\left(\Theta(q)\o(q)\right).
\]


In the following proposition we show that $M$ is a modular form and determine its explicit shape.


\begin{proposition}\label{eq:main-id}
	We have 
	\begin{equation*}
	M(q)
	=
	\frac{4\left(q^2;q^2\right)^2_\infty}{\left(q;q^2\right)^6_\infty}.
	\end{equation*}
	%	Equivalently,
	%	\[
	%	\frac{1}{4q}\bigl(\Theta(q)\omega(q)\bigr)\Big\vert S_{4,1}
	%	=
	%	\prod_{m\ge 1}\frac{(1-q^{8m})^2}{(1-q^{8m-4})^6}.
	%	\]
\end{proposition}


\begin{proof}
%	In \Cref{odissection}, the function $f(q^8)$ contributes only exponents congruent to $0\pmod 4$, while$2q^3\omega(-q^4)$ contributes only exponents congruent to $3\pmod 4$. Hence the terms oresidue classes $1, 2\pmod{4}$ come only from $2q\omega(q)$.
 
%By \Cref{thetadissection}, $\Theta$ is only supported on residue classes $0$ and $1$ modulo $4$. 

We write $\o(q)$ as 
	\[
	\o(q)=\sum_{j=0}^{3}q^j\o_{j}\!\left(q^4\right). 
	\]
\Cref{thetadissection} gives 
\begin{equation}\label{ref2}
M(q)=\Theta(q)\o_1(q)+2\psi\!\left(q^2\right)\o_0(q).
\end{equation}
 By \Cref{odissection}, we have
\[
\o(q)=\frac 12q^{-1}G(q)-\frac 12q^{-1}f\!\left(q^8\right)-\frac 12 q^2\o\!\left(-q^4\right).
\]
Note that only the first term contributes to $\o_0$ and $\o_1$. By \Cref{Fdissection}, we obtain
\[
\frac 12q^{-1}G(q)=\frac 12\sum_{j=0}^3q^{j-1}G_j\!\left(q^4\right).
\]
Then, Lemma 2.7 yields
\begin{align}\label{eqn2}
\o_0(q)&=\frac{G_1(q)}2=\frac{\Theta^2(q)\psi\!\left(q^2\right)}{(q)^2_{\infty}},\quad  \o_1(q)=\frac{G_2(q)}2=\frac{2\Theta(q)\psi^2\!\left(q^2\right)}{(q)^2_{\infty}}.
\end{align}

	%Comparing those residue classes in \Cref{Fdissection}, we obtain
	%\[
	%2q\omega_0\!\left(q^4\right)=qG_1\!\left(q^4\right)\quad \quad \text{and}\quad \quad 2q^2\,\omega_1\left(q^4\right)=q^2G_2\left(q^4\right).
	%\]
%	Therefore, by \Cref{Fdissection}, we have
%	\begin{equation}
%	\omega_0\!\left(q\right)=\frac{G_1\!\left(q^4\right)}{2}=\frac{\Theta^2\!\left(q^4\right)\psi\!\left(q^8\right)}{\left(q^4;q^4\right)_\infty^2},
%	\quad
	%\omega_1\!\left(q\right)=\frac{G_2\!\left(q^4\right)}{2}=\frac{2\Theta\!\left(q^4\right)\psi^2\!\left(q^8\right)}{\left(q^4;q^4\right)^2_\infty}.
	%\end{equation}
By \eqref{ref2}, we then note that
\begin{align*}
&M\!\left(q\right)=%\Theta\!\left(q\right)\omega_1\!\left(q\right)+2\psi\!\left(q^2\right)\omega_0\!\left(q\right)\\
%&\hspace{-0.1 cm}\os{\eqref{eqn2}}=2\Theta\!\left(q\right)\frac{\Theta\!\left(q\right)\psi^2\!\left(q^2\right)}{\left(q\right)^2_\infty}+2\psi\!\left(q^2\right)\frac{\Theta^2\!\left(q\right)\psi\!\left(q^2\right)}{\left(q\right)_\infty^2}=
\frac{4\Theta^2\!\left(q\right)\psi^2\!\left(q^2\right)}{\left(q\right)^2_\infty}%\\
%&\hspace{-0.3 cm}\os{\substack{\eqref{psi}\\(q\mapsto q^2)}}=4\frac{\Theta^2\!\left(q\right)\left(q^{4};q^{4}\right)^4_{\infty}}{\left(q\right)_\infty^2\left(q^2;q^2\right)^2_\infty}%\\
%&\hspace{-0.3 cm}
%\os{\substack{\text{\cite[(5)]{APSY}}}}=4\frac{\left(q^2;q^2\right)^{10}_\infty}{\left(q\right)^6_\infty\left(q^2;q^2\right)^2_\infty}=4\frac{\left(q^2;q^2\right)^{8}_\infty}{\left(q\right)^6_\infty}
=\frac{4\left(q^2;q^2\right)^{2}_\infty}{\left(q;q^2\right)^6_\infty},%\qedhere 
\end{align*}
using \eqref{eqn2}, \eqref{psi}, and \eqref{theta}. This gives the claim.\qedhere 
\end{proof}

We are now ready to show \Cref{ABConj2}.

\begin{proof}[Proof of \Cref{ABConj2}]
Set $D(q):=\frac{(q;q^2)^2_{\infty}}{(-q^2;q^2)_{\infty}}\o(-q)$. Now, by \Cref{lem1},
\begin{equation}\label{B1}
\mathcal{R}\left(S(q)\right)=\mathcal{R}\left(2B(-q)\right)-\mathcal{R}\left(D(q)\right).
\end{equation}
By \Cref{CM}, we have 
\begin{equation}\label{B2}
\mathcal{R}\left(2B(-q)\right)=-2\sum_{n\ge 0}c_B(4n+1)q^n=-\frac{4\left(q^2;q^2\right)^8_{\infty}}{\left(q\right)^7_{\infty}}=-\frac{4\left(q^2;q^2\right)_{\infty}}{\left(q;q^2\right)^7_{\infty}}.
\end{equation}
Next, by \eqref{thetas}, we obtain 
\[
D(q)=\frac{\Theta(-q)\o(-q)}{\left(q^4;q^4\right)_{\infty}}.
\]
 Since, $\Theta(-q)\o(-q)=\sum_{n\ge 0}(-1)^n\d(n)q^n$, we have
\[
\mathcal{R}\left(\Theta(-q)\o(-q)\right)=-M(q).
\]
 As $(q^4;q^4)^{-1}$ contains only powers of $q^{4m}$ ($m\in \mathbb{N}_{ 0}$), we obtain
\[
\mathcal{R}\left(D(q)\right)=-\frac{M(q)}{\left(q\right)_{\infty}}.
\]
By \Cref{eq:main-id}, we have
\[
M(q)=\frac{4\left(q^2;q^2\right)^2_\infty}{\left(q;q^2\right)^6_\infty},
\]
so
\begin{equation}\label{B3}
\mathcal{R}\left(D(q)\right)=-\frac{4\left(q^2;q^2\right)^2_\infty}{\left(q\right)_{\infty}\left(q;q^2\right)^6_\infty}=-\frac{4\left(q^2;q^2\right)_{\infty}}{\left(q;q^2\right)^7_{\infty}}.
\end{equation}
The right-hand sides of \eqref{B2} and \eqref{B3} are equal, hence
\[
\mathcal{R}\left(2B(-q)\right)=\mathcal{R}\left(D(q)\right).
\]
Thus, by \eqref{B1}, we conclude the proof.\qedhere 


%and consequently, it implies that
%\[
%\mathcal{R}\left(D(q)\right)\!=\!\frac{\mathcal{R}\left(\Theta(-q)\o(-q)\right)}{\left(q\right)_{\infty}}\!=\!-\frac{\sum_{n\ge 0}\d(4n+1)q^n}{\left(q\right)_{\infty}}\!=\!-\frac{M(q)}{\left(q\right)_{\infty}}\!=\!-\frac{4\left(q^2;q^2\right)_{\infty}}{\left(q;q^2\right)^7_{\infty}},
%\]
%using \Cref{eq:main-id} in the final step. Plugging this and \eqref{B2} into \eqref{B1}, we conclude the proof of theorem.\qedhere


%Thus to prove \Cref{ABConj2}, by \Cref{lem1}, we need to show that %{\bf \color{blue} K: $c\to d$.} 
%\[
%d(4n+1)=-4\frac{\left(q^2;q^2\right)_{\infty}}{\left(q;q^2\right)^7_{\infty}},%\equiv 0\pmod 4,
%\] 
%where $\sum_{n\ge 0}d(n)q^n:=\frac{(q;q^2)^2_{\infty}}{(-q^2;q^2)_{\infty}}\o(-q)$. To prove this, first note that
%\[
%\frac{\left(q;q^2\right)^2_{\infty}}{\left(-q^2;q^2\right)_{\infty}}\o(-q)=\frac{\left(q\right)_{\infty}\left(q;q^2\right)_{\infty}}{\left(q^4;q^4\right)_{\infty}}\o(-q)=\frac{\Theta(-q)\o(-q)}{\left(q^4;q^4\right)_{\infty}},
%\]
%using \eqref{thetas}. Now $\Theta(-q)\o(-q)=\sum_{n\ge 0}(-1)^n\d(n)q^n$ implies that 
%\[
%\left(\Theta(-q)\o(-q)\right)|S_{4,1}=-\sum_{n\ge 0}\d\left(4n+1\right)q^n=-M(q),
%\]
 %it thus suffices to show that, for $n\in \mathbb{N}_0$, %{\bf KB: Stopped here. \color{blue} K: Explained in details file.} 
%\begin{align*}
%&\textnormal{coeff}_{\left[q^{4n+1}\right]}\left(\left(q;q^2\right)^2_{\infty}\left(q^2;q^2\right)_{\infty}\o(-q)\right)\!=\!
%\textnormal{coeff}_{\left[q^{4n+1}\right]}\left(\left(q\right)_{\infty}\!\left(q;q^2\right)_{\infty}\o(-q)\right)%=\!\textnormal{coeff}_{\left[q^{4n+1}\right]}\left(\frac{\left(q\right)_{\infty}}{(-q)_{\infty}}\o(-q)\right)\nonumber\\
%\!&=\!\textnormal{coeff}_{\left[q^{4n+1}\right]}\left(\o(-q)\Theta(-q)\right)= 0\nonumber\\%\!\pmod 4\nonumber\\
%&\os{\text{\cite[(2.2.12)]{Andbook}}}=\textnormal{coeff}_{\left[q^{4n+1}\right]}\left(\o(-q)\sum_{n\in \mathbb{Z}}(-1)^n q^{n^2}\right)\\
%&\hspace{-0.6 cm}\os{\text{\cite[(2.2.12)]{Andbook}}}=\textnormal{coeff}_{\left[q^{4n+1}\right]}\left(\o(-q)\sum_{n\in \mathbb{Z}}(-1)^n q^{n^2}\right)\os{!}\equiv 0\pmod 4\nonumber\\
%&\Leftrightarrow \textnormal{coeff}_{\left[q^{4n+1}\right]}\left(\Theta(q)\o(q)\right)\equiv 0\pmod 4.
%\end{align*}
%This indeed holds by \Cref{eq:main-id}.\qedhere 
\end{proof}






\section{Open questions}\label{sec5}

  We conclude this paper by proposing following conjectures related to congruences of $c(n)$ for future research.

\begin{conjecture}\label{conj1}
	For every $n\in \mathbb{N}_0$, we have\footnote{We have checked this numerically for $32n+23\le 5000$.}
	\[
	c(32n+23)\equiv 0 \pmod 8.
	\]
\end{conjecture}

 More generally, we propose the following infinite family of congruences for $c(n)$ modulo $4$ and $8$. 

\begin{conjecture}\label{conj2}
For $k,n\in \mathbb{N}_0$, we have
\begin{align*}
c\left(2^{2k+3}n+\frac{11\cdot 4^k+1}{3}\right)\equiv 0\pmod{4},\\
c\left(2^{2k+3}n+\frac{17\cdot 4^k+1}{3}\right)\equiv 0\pmod{8},\\
c\left(2^{2k+4}n+\frac{38\cdot 4^k+1}{3}\right)\equiv 0\pmod{4}.
\end{align*}	
\end{conjecture}

\begin{remark}
{\it Note that we settle the case $k=0$ of \Cref{conj2} in Theorems \ref{ABConj1}--\ref{newthm}. The case $k=1$ (see the congruence modulo $8$) is \Cref{conj1}.}
\end{remark}






%{\bf \color{blue} K: removed all the references not used.}

\begin{thebibliography}{999}
	
\bibitem{A}	G. Andrews, {\it Mordell integrals and Ramanujan’s ‘lost’ notebook’}, pp. 10–48, Analytic number theory, Philadelphia (1980), Lecture Notes in Mathematics 889 (Springer, Berlin, 1981).

	\bibitem{A0} G. Andrews, {\it Ramanujan's ``lost'' notebook I. Partial $\theta$-functions}, Adv. Math. {\bf 41} (1981), 137--172.
	
	%	\bibitem{Andbook} G. Andrews, {\it The theory of partitions}, Cambridge University Press, Cambridge, 1998.
	
%	\bibitem{AB} G. Andrews and B. Berndt, {\it Ramanujan's lost notebook}, Part II, Springer, New York, 2010.
	
%	\bibitem{AB5} G. Andrews and B. Berndt, {\it Ramanujan’s lost notebook}, Part V, Springer, New York, 2018.
	
\bibitem{AB} G. Andrews and M. Bachraoui, {\it Congruences for two-color partitions with odd smallest part},  arXiv:2410.14190v2.

%\bibitem{ADY} G. Andrews, A. Dixit, A. Yee, {\it Partitions associated with the Ramanujan/Watson mock theta functions $\o(q)$, $\nu(q)$ and $\phi(q)$}, Res. Number Theory (2015) {\bf 1}, p. 19.

\bibitem{APSY} G. Andrews, D. Passary, J. Sellers, and A. Yee, {\it Congruences related to the Ramanujan/Watson mock theta functions $\o(q)$ and $\nu(q)$}, Ramanujan J. {\bf 43} (2017), 347--357.

%\bibitem{B} B. Berndt, {\it Ramanujan’s Notebooks}, Part III. Springer, New York (1991).

	\bibitem{BB} B. Berndt and R. Rankin, {\it Ramanujan: letters and commentary}, History of Mathematics {\bf 9} (1995), Amer. Math. Soc., Providence, RI.


\bibitem{CM} S. Chan and R. Mao, {\it Two congruences for Appell--Lerch sums}, Int. J. Number Theory {\bf 8} (2012), 111--123.

%\bibitem{CW} D. Chen and L. Wang, {\it Representations of mock theta functions}, Adv. Math. {\bf 365} (2020), p. 107037.

%\bibitem{G}
%F. Garvan, {\it Universal mock theta functions and two-variable Hecke--Rogers identities},
%Ramanujan J. {\bf 36} (2015), 267--289.


%\bibitem{HM} D. Hickerson and E. Mortenson, {\it Hecke-type double sums, Appell–Lerch sums, and mock theta functions}, Proc. Lond. Math. Soc.  {\bf 109} (2014), 382--422.

%\bibitem{KS} S. Kang and H. Swisher, {\it Mock theta functions of order $2$ and their shadow computations}, Bull. Korean Math. Soc. {\bf 54} (2017), 2155--2163.


\bibitem{Mao} R. Mao, {\it Arithmetic properties of the coefficients of the mock theta function $B(q)$}, Bull. Aust. Math. Soc. {\bf 102} (2020), 50--58.


	
\bibitem{McIntosh} R. McIntosh, {\it Second order mock theta functions}, Canad. Math. Bull. {\bf 50} (2007), 284--290.

%\bibitem{Ono} K. Ono, {\it The web of modularity: arithmetic of the coefficients of modular forms and $q$-series}, CBMS Regional Conference Series in Mathematics {\bf 102} (2003), AMS, Providence, RI.
    

\bibitem{W} G.~Watson, {\it The final problem: An account of the mock-theta functions}, J.~Lond. Math.~Soc. (2) \textbf{11} (1936), 55--80.
   
    
    %\bibitem{Zw} S. Zwegers, {\it Mock theta functions}, Ph.D. Thesis, Universit\"{a}t Utrecht (2002).
    
  
\end{thebibliography}

\end{document}
